% GENERATED FILE. Edit canonical lessons, then run npm run generate:books.
% canonical-source-hash: fnv1a64:93fefb1fc782b780
% canonical-lessons: PA-C01-sat-sri-akal, PA-C01-namaste, PA-C01-han-nahin, PA-W01-ha, PA-W01-aa-matra, PA-W01-bindi, PA-W01-haan-assemble, PA-W01-haan-guided-copy, PA-W01-haan-delayed-copy, PA-W01-haan-dictation, PA-C01-dhanvaad, PA-C01-shukriya, PA-C01-practice

\chapter{Greetings}
\label{ch:greetings}
\hlchaptermodality{1}

\begin{chapteropening}
\textbf{By the end of this chapter:} \emph{I can greet someone in Punjabi and keep reading, listening, independent writing, and speaking separately measurable when I answer yes or no.}

Complete four short, separately scored checks: read \pa{ਹਾਂ}, distinguish heard \pa{ਹਾਂ} and \pa{ਨਹੀਂ}, write \pa{ਹਾਂ} from sound with no visible model, and answer aloud without borrowing credit across skills.
\end{chapteropening}

\section[sat srī akāl]{\emph{sat srī akāl} — the Sikh greeting by ear}

\label{lesson:PA-C01-sat-sri-akal}



\begin{quote}
Begin with one respectful greeting you can hear and return. The printed Gurmukhi is a label today, not a decoding task.
\end{quote}

\subsection*{You'll want to know first — listen and answer}
{[}YOU HEAR: \emph{sat srī akāl}{]}. It serves as both \textbf{hello} and \textbf{goodbye} among Sikhs. Answer by repeating it: {[}YOU SAY: \emph{sat srī akāl}{]}. One careful return is enough.

The phrase is also a small statement of faith: \emph{sat} is “truth,” \emph{srī} is an honorific, and \emph{akāl} is “timeless” — \emph{a-} “not” plus \emph{kāl} “time.” \emph{Akāl} names the Timeless One. \emph{Sat} descends from Sanskrit \emph{satya}, from the same old being-and-truth family remembered by English \textbf{sooth} and \textbf{soothsayer}.

\begin{culture}
This is specifically a Sikh greeting, often said with folded hands. Do not treat it as a generic label for every Punjabi speaker. The next lesson gives a broader pan-Indian greeting.
\end{culture}

\begin{tcolorbox}[breakable,colback=teal!4,colframe=teal!35!black,title={Before you move on}]
Does this greeting serve only as hello? (\textbf{No — hello and goodbye.})
\end{tcolorbox}
\section[namaste]{\emph{namaste} — the shared greeting by ear}

\label{lesson:PA-C01-namaste}



\begin{quote}
Return \emph{sat srī akāl} once. Now meet a broader greeting. Its printed form is still only a label; no sign must be decoded.
\end{quote}

\subsection*{The exchange — listen and answer}
{[}YOU HEAR: \emph{namaste}{]}. It is a respectful \textbf{hello} used across communities and across much of India. Answer in kind: {[}YOU SAY: \emph{namaste}{]}.

Sanskrit \emph{namas} means “a bow” and \emph{te} means “to you”: “a bow to you.” That is why the same word appears in several modern languages. Shared history does not make every community’s greeting interchangeable: use \emph{sat srī akāl} when that Sikh context is present, and \emph{namaste} as the broader greeting taught here.

\begin{tcolorbox}[breakable,colback=teal!4,colframe=teal!35!black,title={Before you move on}]
What do the two Sanskrit pieces mean? (\emph{Namas}, “a bow,” plus \emph{te}, “to you.”)
\end{tcolorbox}
\section[hā̃ / nahī̃]{\emph{hā̃} / \emph{nahī̃} — “yes” and “no” by ear}

\label{lesson:PA-C01-han-nahin}



\begin{quote}
Two tiny answers no conversation survives without. Today they belong to your ear and voice. The printed forms are labels, not a reading test.
\end{quote}

\subsection*{You'll want to know first — meaning before spelling}
{[}YOU HEAR: \emph{hā̃}{]} — \textbf{yes}. Hold the vowel, then let the ending pass gently through the nose.

{[}YOU HEAR: \emph{nahī̃}{]} — \textbf{no}. The second syllable is long, and its ending is nasal too.

Answer aloud once: {[}YOU SAY: \emph{hā̃} — yes{]}. Then once: {[}YOU SAY: \emph{nahī̃} — no{]}. Do not take the Gurmukhi apart yet. The next three tiny sessions will teach one piece at a time before you are asked to read or write \textbf{\pa{ਹਾਂ}}.

\begin{cousinweb}
Punjabi \emph{nahī̃} is built on ancient \textbf{na-}, from the same Proto-Indo-European \textbf{*ne} behind English \textbf{no}, \textbf{not}, and \textbf{none}. That history is a memory hook for the sound and meaning; it is not a spelling cue.
\end{cousinweb}

\begin{tcolorbox}[breakable,colback=teal!4,colframe=teal!35!black,title={Before you move on}]
Say the two answers in order: \textbf{yes, no} — \emph{hā̃, nahī̃}.
\end{tcolorbox}
\section[ha]{\pa{ਹ} — the consonant ha}

\label{lesson:PA-W01-ha}



\begin{quote}
You can already hear and say \emph{hā̃}, “yes.” Today your hand meets only its first consonant. The whole word is still not a reading or writing task.
\end{quote}

\subsection*{Script}
Gurmukhi is an \textbf{abugida}: a consonant already carries a vowel, and you write a vowel mark only when you want a different one.

\begin{quote}
\pa{ਹ}
\end{quote}

That is \textbf{ha} — the consonant \emph{h} with an \textbf{a} already inside it, the \emph{a} of English \emph{about}. Nobody wrote that vowel. A bare consonant is understood to carry it, and it is called the \textbf{inherent vowel}.

Now the name, because it tells you what this script is for. \textbf{Gurmukhi} means \textquotedblleft{}from the mouth of the Guru.\textquotedblright{} It was standardised in the sixteenth century by Guru Angad, the second Sikh Guru, to write down what had until then been recited — and that origin is why it is a \textbf{tidier} script than its neighbours: it was designed in one go, for one purpose, rather than drifting into shape over centuries.

Look at the top. There is a \textbf{head-line} across it, as in Devanagari, and in running text it joins up so a word appears to hang from a line. Gujarati is the sister that erases this line; Gurmukhi keeps it.

\subsection*{Writing — observe and trace}
Keep \textbf{\pa{ਹ}} visible. Follow its printed path once with your finger, then trace one printed model once with a pencil. Say \emph{ha} after your hand stops. Do not hide the model and do not try to write \textbf{\pa{ਹਾਂ}} yet.

\begin{tcolorbox}[breakable,colback=teal!4,colframe=teal!35!black,title={Before you move on}]
What vowel does a bare consonant carry? (\textbf{a} — nobody writes it.) What does \emph{Gurmukhi} mean? (\textbf{From the mouth of the Guru}.) And is the head-line there? (\textbf{Yes} — kept, not erased.)

Source: \href{https://www.unicode.org/charts/PDF/U0A00.pdf}{Unicode Gurmukhi chart}.
\end{tcolorbox}
\section[-ā]{\pa{ਾ} — the vowel sign that replaces the inherent a}

\label{lesson:PA-W01-aa-matra}



\begin{quote}
Write the piece from the previous lesson before adding another.
\end{quote}

\subsection*{Script}
To get a different vowel you \textbf{replace} the inherent one, with a mark called a \textbf{mātrā}.

\begin{quote}
\pa{ਾ}
\end{quote}

A single upright stroke standing to the \textbf{right} of its consonant, dropping from the head-line. Attached, it swaps the short \emph{a} for \textbf{ā}, the \emph{a} of English \emph{father}.

\begin{quote}
\textbf{\pa{ਹ}} + \textbf{\pa{ਾ}} $\to$ \textbf{\pa{ਹਾ}}
\end{quote}

\emph{hā.} One beat — a mātrā changes the vowel, it never adds a beat.

You are one mark away from a word.

\subsection*{Writing — observe and trace}
Keep \textbf{\pa{ਾ}} visible beside \textbf{\pa{ਹ}}. Finger-trace the vowel sign once, then trace one printed model once with a pencil. Point to \textbf{\pa{ਹ}}, then \textbf{\pa{ਾ}}, and say \emph{hā}. Keep the model visible; the nasal ending has not been taught yet.

\begin{tcolorbox}[breakable,colback=teal!4,colframe=teal!35!black,title={Before you move on}]
What does a mātrā do to the inherent vowel? (\textbf{Replaces} it.) Where does this one sit? (To the \textbf{right}.) And does it add a beat? (\textbf{No}.)

Source: \href{https://www.unicode.org/charts/PDF/U0A00.pdf}{Unicode Gurmukhi chart}.
\end{tcolorbox}
\section[-ṁ]{\pa{ਂ} — a dot above the line that nasalises the vowel}

\label{lesson:PA-W01-bindi}



\begin{quote}
Write the piece from the previous lesson before adding another.
\end{quote}

\subsection*{Script}
Here is a mark that is not a vowel at all.

\begin{quote}
\pa{ਂ}
\end{quote}

This is the \textbf{bindi}, a dot sitting \textbf{above the head-line}. It does not change which vowel you say. It says: \textbf{send that vowel through your nose.}

English does this constantly and never writes it. Say \emph{bad}, then say \emph{ban} — the vowel in the second is nasalised before you ever reach the \emph{n}. Punjabi writes what English leaves implicit.

\begin{quote}
\textbf{\pa{ਹਾ}} + \textbf{\pa{ਂ}} $\to$ \textbf{\pa{ਹਾਂ}}
\end{quote}

\textbf{hāṁ} — and that is \textbf{yes}, the word you have been saying since your first lesson.

\textbf{Nasalisation is meaning-bearing here.} It is not an accent or a flourish; drop the dot and you have a different word. That is why the script spends a character on it.

\subsection*{Writing — observe and trace}
Keep \textbf{\pa{ਂ}} visible above the head-line. Point to it once, finger-trace the dot once, then trace one printed model with a pencil. Say “nasal” as your hand stops. Do not copy the whole word yet; assembly is the next separate session.

\begin{tcolorbox}[breakable,colback=teal!4,colframe=teal!35!black,title={Before you move on}]
What does the bindi do? (\textbf{Nasalises} the vowel — sends it through the nose.) Does it change which vowel you say? (\textbf{No}.) And what word did it complete? (\textbf{\pa{ਹਾਂ}} — \emph{hāṁ}, \textquotedblleft{}yes\textquotedblright{}.)

Source: \href{https://www.unicode.org/charts/PDF/U0A00.pdf}{Unicode Gurmukhi chart}.
\end{tcolorbox}
\section[hā̃]{\pa{ਹਾਂ} — three known pieces become one known answer}

\label{lesson:PA-W01-haan-assemble}



\begin{quote}
Point to the three pieces you met separately: \textbf{\pa{ਹ}}, \textbf{\pa{ਾ}}, \textbf{\pa{ਂ}}. Say their jobs: \emph{h}, long \emph{ā}, nasal ending.
\end{quote}

\subsection*{Script — assemble with no new sign}
Slide the known pieces together from left to right:

\begin{quote}
\textbf{\pa{ਹ}} + \textbf{\pa{ਾ}} + \textbf{\pa{ਂ}} $\to$ \textbf{\pa{ਹਾਂ}}
\end{quote}

Read \emph{hā̃}, “yes.” Point once to each piece while you say its part, then read the whole answer once at natural speed. Your pencil stays down today; assembly and copying are separate steps.

\begin{tcolorbox}[breakable,colback=teal!4,colframe=teal!35!black,title={Before you move on}]
Which piece nasalises the ending? (\textbf{\pa{ਂ}}, the bindi.)
\end{tcolorbox}
\section[hā̃]{\pa{ਹਾਂ} — one supported copy}

\label{lesson:PA-W01-haan-guided-copy}



\begin{quote}
Keep \textbf{\pa{ਹਾਂ}} visible. Point to \textbf{\pa{ਹ}}, then attached \textbf{\pa{ਾ}}, then the bindi \textbf{\pa{ਂ}}. Read the known answer \emph{hā̃}, “yes.”
\end{quote}

\subsection*{Writing — guided copy}
Copy \textbf{\pa{ਹਾਂ}} once beside the model. Look back after every piece if you need to. Touch the three pieces in your copy and read \emph{hā̃}. Large and slow is correct; memory is not being tested yet.

\begin{tcolorbox}[breakable,colback=teal!4,colframe=teal!35!black,title={Before you move on}]
Compare one piece at a time. Repair one piece if needed, then stop.
\end{tcolorbox}
\section[hā̃]{\pa{ਹਾਂ} — look, hide, write, compare}

\label{lesson:PA-W01-haan-delayed-copy}



\begin{quote}
Look at \textbf{\pa{ਹਾਂ}} for five seconds. Point to its three pieces. Do not copy it yet.
\end{quote}

\subsection*{Writing — delayed copy}
1. cover \textbf{\pa{ਹਾਂ}} 2. wait five seconds 3. write it once from memory 4. uncover the model and compare \textbf{\pa{ਹ}}, \textbf{\pa{ਾ}}, then \textbf{\pa{ਂ}}

Circle and repair only a differing piece. Repair is part of the exercise, not a failure.

\begin{tcolorbox}[breakable,colback=teal!4,colframe=teal!35!black,title={Before you move on}]
Read your repaired word as \emph{hā̃}. Stop after one memory attempt.
\end{tcolorbox}
\section[hā̃]{One heard answer — write what you know}

\label{lesson:PA-W01-haan-dictation}



\begin{quote}
Cover the answer lower on this page. You will hear one familiar answer. No new Punjabi or new sign is hiding here.
\end{quote}

\subsection*{Writing — short dictation}
{[}YOU HEAR: \emph{hā̃} — “yes”{]}

Write the Punjabi word from sound alone. Now uncover and compare: \textbf{\pa{ਹਾਂ}}. Check \textbf{\pa{ਹ}} first, attached \textbf{\pa{ਾ}} second, and \textbf{\pa{ਂ}} above the line last. Repair one place if needed.

\begin{tcolorbox}[breakable,colback=teal!4,colframe=teal!35!black,title={Before you move on}]
Writing from sound completes this tiny runway. Later lessons must grow from one known word to phrases, messages, and timed exam tasks.
\end{tcolorbox}
\section[dhannavād]{\emph{dhannavād} — “thank you,” the Sanskrit way}

\label{lesson:PA-C01-dhanvaad}



\begin{quote}
You can greet and answer yes or no. Add one spoken thanks. The Gurmukhi above is a label, not a word you must decode.
\end{quote}

\subsection*{You'll want to know first — listen and use}
{[}YOU HEAR: \emph{dhannavād}{]} — \textbf{thank you}. Return it once: {[}YOU SAY: \emph{dhannavād}{]}.

It descends from Sanskrit \emph{dhanya}, “worthy” or “blessed,” plus \emph{vāda}, “a saying.” This is Punjabi’s Sanskritic thanks. The next session adds a Perso-Arabic choice with the same everyday function.

\begin{tcolorbox}[breakable,colback=teal!4,colframe=teal!35!black,title={Before you move on}]
Is this the Sanskritic or Perso-Arabic choice? (\textbf{Sanskritic.})
\end{tcolorbox}
\section[shukrīā]{\emph{shukrīā} — “thank you,” the Persian way}

\label{lesson:PA-C01-shukriya}



\begin{quote}
Say the Sanskritic thanks \emph{dhannavād}. Now add a second spoken choice. No new Gurmukhi sign is required.
\end{quote}

\subsection*{You'll want to know first — listen and use}
{[}YOU HEAR: \emph{shukrīā}{]} — \textbf{thank you}. Return it once: {[}YOU SAY: \emph{shukrīā}{]}.

The word came through Persian from Arabic \emph{shukr}, “gratitude,” the same root heard in Arabic \emph{shukran}. Punjabi carries both a Sanskritic thanks, \emph{dhannavād}, and this Perso-Arabic one. Learn the contrast by history and sound; the dotted borrowed-sound spelling can wait for its own sign lesson.

\begin{tcolorbox}[breakable,colback=teal!4,colframe=teal!35!black,title={Before you move on}]
Say both choices in historical order: Sanskritic \emph{dhannavād}, then Perso-Arabic \emph{shukrīā}.
\end{tcolorbox}
\section[Practice]{Chapter 1 — four small checks, four separate scores}

\label{lesson:PA-C01-practice}



\begin{quote}
No new Punjabi. Complete each tiny check on its own. A correct answer in one skill never erases a miss in another.
\end{quote}

\subsection*{How to answer — greetings and courtesy first}
{[}YOU SAY: \emph{sat srī akāl}{]}, {[}YOU SAY: \emph{namaste}{]}, {[}YOU SAY: \emph{dhannavād}{]}, {[}YOU SAY: \emph{shukrīā}{]}. Record this courtesy check separately from the yes/no speaking check below.

\subsection*{Script — see and understand}
Read \textbf{\pa{ਹਾਂ}} without saying the answer from memory first. Record only the reading result.

\subsection*{You'll want to know first — hear and choose}
{[}YOU HEAR: \emph{nahī̃}{]} Choose the meaning without looking at a written Punjabi answer. Record only the listening result.

\subsection*{Writing — one independent answer}
Cover every Gurmukhi answer on this page. Hear \emph{hā̃}, “yes,” once and write it from sound alone. Do not award writing credit for a romanized answer or a copy. After scoring, uncover \textbf{\pa{ਹਾਂ}} and repair one piece if needed.

\subsection*{Your turn — answer aloud}
{[}YOU SAY: \emph{hā̃} — yes{]}. Record only the spoken response. The written answer from the previous check cannot supply speaking credit.

\begin{tcolorbox}[breakable,colback=teal!4,colframe=teal!35!black,title={Before you move on}]
Name the four separately scored skills: reading, listening, writing, and speaking. Repeat only the skill that missed; do not average it away.
\end{tcolorbox}
